\documentclass[letterpaper]{article}
\usepackage[submission]{aaai2026}
\usepackage{times}
\usepackage{helvet}
\usepackage{courier}
\usepackage{graphicx}
\usepackage{caption}
\usepackage{natbib}
\usepackage{url}
\usepackage{booktabs}
\usepackage{amsmath}

\title{Does Shared Investigation State Improve Human--AI Coordination\\on Reverse Engineering Tasks?}

\author{
    Anonymous Author(s)
}

\begin{document}
\maketitle

\begin{abstract}
We study whether making an LLM agent's investigation state explicit and shared---rather than implicit in conversation history---improves reasoning quality and responsiveness to human guidance on reverse engineering tasks.
We introduce a minimal agent framework in which an LLM iteratively analyzes Linux ELF binaries by forming hypotheses, calling tools, and updating a structured investigation state.
Two agent configurations share the same loop, tools, and interventions; they differ only in whether the investigation state is visible to the model at each turn.
Scripted interventions, derived from independently authored human writeups via a reproducible LLM-assisted pipeline, serve as proxies for human collaborator input.
In preliminary experiments on crackme challenges using DeepSeek R1 14B, the structured agent solved 67\% of trials while the baseline agent solved none, despite receiving identical interventions.
The baseline agent exhibited a pronounced \emph{action gap}---describing intended analyses without executing them---and failed to act on interventions delivered through conversation history alone.
These results suggest that explicit shared state may be a necessary coordination substrate for effective human--AI teaming on investigative tasks, particularly with smaller language models.
\end{abstract}

\section{Introduction}

Reverse engineering is an inherently iterative, hypothesis-driven process.
An analyst inspects a binary, forms hypotheses about its behavior, tests them through controlled experiments, refines understanding, and repeats.
This investigative structure makes reverse engineering a compelling testbed for studying human--AI collaboration: the task requires sustained multi-step reasoning, strategic tool use, and the ability to incorporate external guidance mid-investigation.

Recent work has demonstrated that large language models can perform aspects of reverse engineering when given appropriate tools \citep{shao2024nyx, yang2024rag}.
However, most agent frameworks maintain investigation context implicitly through conversation history.
As the investigation progresses and context grows, the model must reconstruct its current understanding from an increasingly long and unstructured message log.
This raises a question central to human--AI teaming:

\begin{quote}
\emph{Does making an LLM agent's investigation state explicit and shared---rather than implicit in conversation history---improve its reasoning quality, tool usage, and responsiveness to human guidance?}
\end{quote}

We address this question through a controlled experiment comparing two agent configurations on the same reverse engineering tasks.
Both agents share an identical loop, tool set, and intervention schedule.
The only difference is \emph{prompt construction}: the baseline agent sees only its conversation history, while the structured agent additionally receives a rendered investigation state---observations, hypotheses, actions taken, and current focus---at every turn.

Our contributions are:
\begin{enumerate}
    \item A minimal agent framework for reverse engineering that supports controlled comparison of state visibility and intervention strength.
    \item A reproducible pipeline for deriving tiered interventions from human writeups, serving as proxies for human collaborator input.
    \item Preliminary experimental evidence that explicit shared state substantially improves tool usage, intervention uptake, and solve rate, even with small (14B parameter) language models.
\end{enumerate}

\section{Related Work}

\paragraph{LLM agents for security tasks.}
Several recent systems apply LLM agents to cybersecurity challenges.
NYX \citep{shao2024nyx} uses an LLM with tool access to solve CTF challenges across multiple categories.
\citet{yang2024rag} augment an LLM agent with retrieval over security documentation.
These systems demonstrate that LLMs can perform multi-step security analysis, but they do not study the effect of state representation on reasoning quality.

\paragraph{Structured reasoning in LLM agents.}
Chain-of-thought prompting \citep{wei2022chain} and its variants improve LLM reasoning by eliciting intermediate steps.
ReAct \citep{yao2023react} interleaves reasoning and action in a shared trace.
Reflexion \citep{shinn2023reflexion} adds explicit self-evaluation.
Our work differs in that we externalize the \emph{investigation state} as a shared object between the agent and a human collaborator, rather than using it solely as an internal reasoning aid.

\paragraph{Human--AI teaming.}
Research on human--AI collaboration has studied shared mental models \citep{johnson2014coactive}, adjustable autonomy \citep{bradshaw2004making}, and mixed-initiative interaction \citep{horvitz1999principles}.
Our intervention mechanism---typed state mutations rather than free-text hints---is inspired by the concept of \emph{common ground} in collaborative problem-solving \citep{clark1996using}: the shared investigation state serves as a coordination substrate that both parties can read and write.

\section{System Design}

\subsection{Agent Loop}

Both agent configurations execute the same loop (Algorithm~\ref{alg:loop}).
At each step, the agent receives a prompt, generates a response, and the system parses the response for: (1)~a tool call, (2)~a structured state update, and (3)~a solution declaration.
If a tool is called, the result is fed back as the next user message.
Interventions are checked after each step via a stall-based trigger mechanism.

\begin{figure}[t]
\centering
\small
\begin{tabular}{l}
\toprule
\textbf{Algorithm 1:} Agent Investigation Loop \\
\midrule
\textbf{Input:} challenge $C$, config (structured, intervention mode) \\
\textbf{Output:} InvestigationState $S$ \\
\\
$S \leftarrow$ \textsc{InitState}($C$) \\
$H \leftarrow []$ \hfill \textit{// message history} \\
$Q \leftarrow$ \textsc{LoadInterventions}($C$, mode) \\
\\
\textbf{for} step $= 1$ \textbf{to} max\_steps \textbf{do} \\
\quad $P \leftarrow$ \textsc{BuildPrompt}($S$, $H$, structured) \\
\quad $R \leftarrow$ \textsc{LLM}($P$) \\
\quad \textbf{if} \textsc{HasToolCall}($R$) \textbf{then} \\
\quad\quad $output \leftarrow$ \textsc{Execute}(\textsc{ParseTool}($R$)) \\
\quad\quad $S \leftarrow$ \textsc{RecordAction}($S$, tool, output) \\
\quad \textbf{if} structured \textbf{then} \\
\quad\quad $S \leftarrow$ \textsc{ApplyStateUpdate}($S$, \textsc{ParseJSON}($R$)) \\
\quad \textbf{if} \textsc{HasSolution}($R$) or \textsc{BinaryAccepted}(output) \textbf{then} \\
\quad\quad \textbf{return} $S$ with solved $= $ true \\
\quad $iv \leftarrow Q$.\textsc{CheckStall}(step, new\_tool\_used) \\
\quad \textbf{if} $iv \neq$ null \textbf{then} \\
\quad\quad $S \leftarrow$ \textsc{ApplyIntervention}($S$, $iv$) \\
\quad $H \leftarrow H + [R, output]$ \\
\textbf{return} $S$ \\
\bottomrule
\end{tabular}
\caption{The shared agent loop. Both baseline and structured variants execute identical logic; only \textsc{BuildPrompt} differs.}
\label{alg:loop}
\end{figure}

\subsection{Experimental Conditions}

We use a factorial design crossing state visibility with intervention strength:

\begin{table}[t]
\centering
\caption{Experimental conditions. All conditions use the same loop, tools, and challenge set.}
\label{tab:conditions}
\begin{tabular}{clcc}
\toprule
 & Condition & State Visible & Interventions \\
\midrule
A & Baseline & No & None \\
B & Baseline & No & Light \\
C & Structured & Yes & None \\
D & Structured & Yes & Light \\
E & Structured & Yes & Strong \\
\bottomrule
\end{tabular}
\end{table}

\paragraph{State visibility.}
In the \emph{baseline} condition, the model receives only a system prompt and its conversation history.
The \texttt{InvestigationState} object is maintained internally by the loop but is never included in the prompt.
In the \emph{structured} condition, the state is rendered as human-readable text and prepended to every user message.
The model is additionally instructed to emit a JSON state update block at the end of each response.

\paragraph{Intervention strength.}
Light interventions provide directional guidance: \texttt{add\_hypothesis}, \texttt{reject\_hypothesis}, \texttt{mark\_observation\_important}.
Strong interventions prescribe actions: \texttt{suggest\_action}, \texttt{reprioritize\_next\_steps}.
Interventions are delivered via a stall-based trigger: after $k$ consecutive steps with no new tool called, the next intervention from a tiered queue is applied.

\subsection{Investigation State}

The shared state is a structured object containing:

\begin{itemize}
    \item \textbf{Observations}: facts discovered via tools, with source attribution.
    \item \textbf{Hypotheses}: testable claims with status (active/supported/refuted), confidence level, and \emph{origin} (agent or intervention).
    \item \textbf{Actions}: tool calls and their results.
    \item \textbf{Current focus}: what the agent is currently investigating.
    \item \textbf{Current understanding}: a free-text summary the agent maintains.
\end{itemize}

The \texttt{origin} field on hypotheses is the primary mechanism for measuring intervention uptake: it tracks whether a hypothesis was introduced by the agent or injected by an intervention.

\subsection{Tools}

The agent has access to four tools: \texttt{file()} (binary metadata), \texttt{strings()} (printable string extraction), \texttt{run\_binary(input)} (execute the binary with a given argument), and \texttt{python\_eval(code)} (run a Python snippet).
Static analysis tools execute locally; \texttt{run\_binary} executes x86-64 Linux ELF binaries via Docker with Rosetta emulation on ARM hardware.

Tool calls are parsed from the model's natural-language response (e.g., ``\texttt{TOOL: run\_binary(test)}'') rather than through a structured tool-calling API.
This design choice reflects a practical constraint: many open-weight models do not reliably support formal tool-calling schemas, particularly at smaller parameter counts.

\section{Intervention Derivation}

\subsection{Human Reasoning Proxies}

Rather than recruiting human participants for initial experiments, we use independently authored solution writeups as proxies for human reasoning.
For each challenge, we collect multiple writeups from public repositories (e.g., crackmes.one).
After quality filtering (excluding code-only dumps and excessively brief submissions), each writeup is processed through a two-stage LLM-assisted pipeline.

\subsection{Stage 1: Reasoning Step Extraction}

Each writeup is converted into a structured sequence of typed steps: \emph{observation} (a discovered fact), \emph{hypothesis} (a testable claim), \emph{action} (a tool or analysis step taken), and \emph{result} (an outcome).
Extraction is performed by Claude Opus 4 using a constrained prompt that requires direct textual evidence for each step (see Appendix for full prompts).
A normalization pass merges duplicates, removes vague claims, and ensures atomicity.

\subsection{Stage 2: Intervention Synthesis}

All extracted traces for a challenge are provided to the model, which identifies recurring insights across writeups.
Consensus observations (present in $\geq$3 writeups) become intervention candidates.
Interventions are tiered:

\begin{itemize}
    \item \textbf{Tier 1 (light)}: directional---points at \emph{what} to investigate (e.g., ``The binary has a \texttt{checkSerial} function that validates input'').
    \item \textbf{Tier 2 (strong)}: structural---describes the \emph{pattern} without revealing a specific valid input (e.g., ``The validation compares pairs of adjacent characters using an arithmetic relationship'').
\end{itemize}

\paragraph{Human verification.}
All LLM-generated artifacts (extracted steps and synthesized interventions) undergo manual review to ensure: (1)~no fabricated steps, (2)~accurate evidence attribution, (3)~correct type assignments, and (4)~no solution leakage in interventions.
This human-in-the-loop validation is essential because the extraction model may hallucinate reasoning steps not present in the source text or inadvertently encode solution-revealing information in intervention phrasing.

\section{Experimental Setup}

\subsection{Challenge}

We report preliminary results on \texttt{yurisimplekeygen}, a beginner-level crackme from crackmes.one (difficulty 1.5/6, quality 5.0/5).
The challenge is a 64-bit x86-64 ELF binary that validates a 16-character serial key.
Seven independently authored writeups were collected and processed through the intervention derivation pipeline.

\subsection{Model}

All experiments use DeepSeek R1 14B \citep{deepseekai2025deepseekr1}, served via Ollama on a remote GPU server.
DeepSeek R1 produces explicit chain-of-thought reasoning in \texttt{<think>} blocks, which the system captures as a separate trace for analysis.
The 14B parameter count represents a practical constraint: it is the largest model that runs at interactive speeds on available hardware.

\subsection{Protocol}

Each condition is run 3 times with a 15-step budget.
The stall threshold for intervention delivery is $k=3$ consecutive steps with no new tool.
Runs are logged as structured JSONL records capturing per-step tool usage, intervention application, response length, and derived metrics.

\section{Preliminary Results}

\begin{table}[t]
\centering
\caption{Preliminary results on \texttt{yurisimplekeygen} (3 runs per condition, 15-step budget, DeepSeek R1 14B). Conditions B and D shown (both with light interventions).}
\label{tab:results}
\begin{tabular}{lcccc}
\toprule
 & Solve & Tool & Unique & \texttt{run\_binary} \\
Condition & Rate & Calls & Tools & Calls \\
\midrule
B: Baseline + Light & 0/3 & 4.3 & 1.3 & 0.0 \\
D: Structured + Light & 2/3 & 3.0 & 3.3 & 1.7 \\
\bottomrule
\end{tabular}
\end{table}

Table~\ref{tab:results} reports preliminary results.
Three findings emerge:

\paragraph{1. Structured state enables progressive tool use.}
The structured agent followed purposeful investigation arcs (e.g., \texttt{file} $\rightarrow$ \texttt{strings} $\rightarrow$ \texttt{run\_binary}), using 3.3 unique tools on average.
The baseline agent called \texttt{file()} repeatedly (up to 5 times in a single run) without progressing to other tools, averaging only 1.3 unique tools.
Critically, the baseline agent \emph{never} called \texttt{run\_binary} in any run---meaning it never tested any input against the binary.

\paragraph{2. Interventions require state context to be effective.}
Both conditions received identical interventions.
The structured agent acted on interventions within 1--2 steps in runs where it was making tool calls.
The baseline agent completely ignored interventions in all 3 runs, despite receiving the same information in its conversation history.
This suggests that interventions delivered as text in a growing conversation are insufficient; the structured state provides the context needed to make interventions actionable.

\paragraph{3. Baseline exhibits a pronounced action gap.}
The baseline agent produced responses approximately 3$\times$ longer than the structured agent, predominantly consisting of multi-paragraph descriptions of \emph{intended} analyses rather than actual tool calls.
We term this the \emph{action gap}: the model describes what it would do without executing it.
The structured agent, with its visible record of actions already taken and a defined current focus, exhibited this behavior far less frequently.

\section{Limitations}

These results are preliminary.
The sample size (3 runs per condition, 1 challenge) is insufficient for statistical significance.
The current interventions, while improved from earlier iterations, still provide substantial guidance; weaker interventions would better test the state visibility effect in isolation.
DeepSeek R1 14B is a relatively small model that struggles with instruction following; larger models may reduce the baseline's tool-calling failures, potentially narrowing the gap.
We have not yet tested conditions A, C, and E, which are needed to fully disentangle the effects of state visibility and intervention presence.

\section{Conclusion and Future Work}

We presented a controlled framework for studying whether explicit shared investigation state improves human--AI coordination on reverse engineering tasks.
Preliminary results suggest that state visibility is not merely helpful but may be \emph{necessary} for effective tool use and intervention uptake, particularly with smaller language models.

Planned extensions include: (1)~completing the factorial design across all five conditions, (2)~adding 2--3 additional challenges of varying difficulty, (3)~testing with larger models (DeepSeek R1 70B) and frontier models (Claude) to assess whether the effect persists at scale, and (4)~replacing scripted interventions with live human collaborators to validate the proxy-based findings.

\bibliographystyle{aaai2026}
\bibliography{references}

\end{document}
